\section{From the 64 residues back to the arithmetic witness}
\label{sec:arithmetic-transport}

A pair of residues at distance \(s\) can be a geometrically allowed edge while neither endpoint is retained on the host. An odd retained point can incur a charge while the fallback at its orbit receives no credit. Both distinctions matter when interpreting a percentage. This section states the hypotheses under which a quotient profile has an arithmetic interpretation, and derives the two weighted sums whose ratio the experiment computes. It uses the original conductor-one scales and the comparison results established in the preceding studies.

\subsection{What is fixed at the moment of selection}

Fix \(0<\varepsilon<1/52\). In the principal conductor-one specialization, let \(D\) be a raw host at a sufficiently late family point \(x\), with its field, prime-pair row, and point labels. The relevant cyclic group \(G\) has dyadic order \(n\asymp x^h\), where \(h=2+8\varepsilon\). A prime \(p\) has an arithmetic Frobenius step \(s_p\) before the set of retained tops is chosen. The source represents this step using an existence choice and the projection of an arithmetic Frobenius class. Its orbit is based at the least \emph{numeric host label} in the coset. The source also supplies the coordinate indexing the cycle from that base. Among generators of the same even cyclic subgroup, these indices differ by multiplication by an odd unit, which preserves parity. The support decision can therefore depend on one preferred child per node of the quotient tree. An arbitrary preference tree ranges over numerical orders with these minima; a worst-case tree used as an upper certificate is an abstract labeling unless it has separately been realized by an arithmetic host.
\href{https://github.com/nousentllc/riemann-hypothesis/blob/40e56bab026552e089d7afcfeff08d2defb9c005/src/Development/AbelianVarietyTheory/Definitions/GlobalTwisting/ThinScalesIffHostBound.lean#L19-L41}{Frobenius choice and eligible-prime predicate}.
\href{https://github.com/nousentllc/riemann-hypothesis/blob/40e56bab026552e089d7afcfeff08d2defb9c005/src/Formalism/LocalToGlobal/GlobalTheory/IsResidueClassZero.lean#L264-L275}{Canonical base}.

The original call obtains \(\exists\,T\), then builds the record \(\{D\ \mathrm{with}\ TT:=T\}\). It exports a size band, a bound on bad-orbit mass, and a share condition. One can test a stronger requirement than these three public conditions, because the selection proof names the scalar score used to produce the witness. Put \(\delta=x^{-\varepsilon/8}\), let \(\mathcal P\) consist of primes \(p< H=\exp(x^{\varepsilon/3})\) in the source's residue-class-zero category with \(\operatorname{ord}(s_p)\ge x^{\varepsilon/4}\), and write \(\omega_p=\log(p)/p\). Define
\begin{align*}
 B(T)&=\sum_{p\in\mathcal P}\omega_p
   \sum_{y\in T}\mathbf 1_{\{\text{the orbit through }y\text{ fails its windows at }p\}},\\
 M&=C_bx^{-2}\delta n\sum_{p\in\mathcal P}\omega_p,\qquad
 S(T)=\frac{(|T|-\delta n)^2}{2\delta n}+\frac{B(T)}{4M}.
\end{align*}
If \(M=0\), the bad-prime term is zero. The subgroup used for the conductor-one shares is the entire group, so only one cardinality appears in the centered count. In the original product law, each site has the same marginal \(\delta\). The proof obtains a set of positive product mass satisfying \(S(T)\le\mathbb E_\delta S\). The binomial variance gives \(\mathbb E_\delta S\ge(1-\delta)/2\). This lower bound gives a way to exhibit other admissible sets without evaluating the original expectation. Every fixed subset has positive probability under the original law once \(0<\delta<1\).
\href{https://github.com/nousentllc/riemann-hypothesis/blob/40e56bab026552e089d7afcfeff08d2defb9c005/src/Development/AutomorphicArithmetic/ThickFrames/select.lean#L136-L173}{The score in the averaging proof}.
\href{https://github.com/nousentllc/riemann-hypothesis/blob/40e56bab026552e089d7afcfeff08d2defb9c005/src/Development/AutomorphicArithmetic/ThickFrames/topsGroup.lean#L88-L123}{Choice of the original product measure}.
\href{https://github.com/nousentllc/riemann-hypothesis/blob/40e56bab026552e089d7afcfeff08d2defb9c005/src/Development/AutomorphicArithmetic/ShortNorms/hostOfFieldOf.lean#L217-L229}{Insertion of the witness into the host}.

\subsection{Why a nonuniform profile can pass that same score}

Let \(q=64\), and choose one rate \(a_t\ge0\) per residue with \(\sum_t a_t=q\). In a class of size \(n/q\) choose uniformly \(k_t\) points, independently across classes. Round the desired counts \(a_t\delta n/q\) so that \(\sum k_t=N\) is the nearest integer to \(\delta n\). Taking lower integers and filling the remaining slots in order of fractional part gives an error at most one per occupied class and preserves the zeros. Thus each realized mask has exactly the target total to within one half. To use the source's coefficient \(C_b\), which was fixed before these sets were chosen, we must also recover the scale \(o(x^{-2})\) in the bad probability \emph{conditional on a retained point}. Convergence of the unconditioned failure probability to zero would not settle this comparison.

Let \(y\) be a retained anchor and fix a prescribed pattern at \(\ell\) further sites in one class. With population \(m\) and sample \(k\) after conditioning on \(y\), the probability of a specified pattern of \(i\) successes is
\[
 \frac{(k)_i(m-k)_{\ell-i}}{(m)_\ell}
 \ \le\ \frac{m^\ell}{(m)_\ell}
       \left(\frac{k}{m}\right)^i
       \left(1-\frac{k}{m}\right)^{\ell-i}.
\]
Here \((m)_\ell\) denotes a falling factorial. The first factor has logarithm \(O(\ell^2/m)\) for \(\ell\le m/2\). Across all \(q\) classes \(m\ge n/q-1\); the total number inspected in one window or block is \(O(\delta^{-1}\log x)\) for fixed \(\varepsilon\). The comparison factors therefore multiply to \(1+o(1)\), uniformly in the pattern, location, and anchored class. Summing over exactly those patterns that constitute a failure yields a \emph{relative} upper comparison, valid even for a rare event. An additive approximation of order \(\ell^2/m\) would leave that crucial scale to be proved again.

On an eligible orbit write \(L=\operatorname{ord}(s_p)\) and \(u=\lceil1/\delta\rceil\). In an independent comparison cell of \(u\) points let \(\mu\) be its limiting mean count and \(\nu\) its limiting even-position count. A projected orbit visits a coset of the 64 residues periodically. The site probabilities have maximum \(O(\delta)\); their squares sum to \(O(\delta)\). Expanding the logarithm of the empty-cell probability therefore gives the limit \(e^{-\mu}\), uniformly in the starting position. The even-singleton probability is the empty probability times the sum of the odds \(p_i/(1-p_i)\) over the even sites. Its limit is \(\nu e^{-\mu}\). At a homogeneous residue and below it one has \(\nu=\mu/2\). At a coarser node the relevant even sites are defined by its least label. The source tile consists of \(3r+1\) empty cells alternating with \(3r\) even-singleton cells in prescribed order, where \(r\) is fixed once \(\varepsilon\) is fixed. Thus its limiting probability is
\[
 \left(e^{-\mu}\nu^{\,3r/(6r+1)}\right)^{6r+1}.
\]
Suppose that on every occupied type the factor in parentheses strictly exceeds \(1/10\). For a fixed finite menu, choose \(c>1/10\) below all of these factors; a sufficiently late tile succeeds under the comparison law with probability at least \(c^{6r+1}\). Uniformity here concerns the finitely many types and every starting cell, with the original order of limits \(x\to\infty\) at fixed \(\varepsilon\). It supplies no practical first value of \(x\).
\href{https://github.com/nousentllc/riemann-hypothesis/blob/40e56bab026552e089d7afcfeff08d2defb9c005/src/Formalism/SubspaceCounting/GoodOrbitIntervalTileCondition.lean#L9-L26}{Cell parities and three orbit conditions}.
\href{https://github.com/nousentllc/riemann-hypothesis/blob/40e56bab026552e089d7afcfeff08d2defb9c005/src/Formalism/SharpRestriction/GrowthBounds/UnramifiedQuadraticSetting.lean#L237-L251}{Cell scale and window lengths}.

The source searches \(M_0+1\) distinct consecutive trial tiles from each start, with \(M_0\ge(h+3)10^{6r+1}\log x\). The orbit-length hypotheses leave room for the block. Ignore the one tile containing the retained anchor, if present. The remaining disjoint trials are independent in the comparison law; any given start fails with probability at most \((1+o(1))\exp(-M_0 c^{6r+1})\). Multiplication by the \(O(x^h)\) possible starts leaves \(o(x^{-2})\). Applying the relative comparison to each block \emph{before} taking this union is important, since a whole orbit can be far larger than an individual inspection.
\href{https://github.com/nousentllc/riemann-hypothesis/blob/40e56bab026552e089d7afcfeff08d2defb9c005/src/Development/AutomorphicArithmetic/ThickFrames/orbitTerms.lean#L107-L136}{Eligible orbit and trial-block size}.
\href{https://github.com/nousentllc/riemann-hypothesis/blob/40e56bab026552e089d7afcfeff08d2defb9c005/src/Development/AutomorphicArithmetic/ThickFrames/TilesLe.lean#L43-L67}{Delete the trial containing the anchor}.

For completeness, take a common eventual marginal bound \(A_*\delta\), with \(A_*<3\), encompassing the rate profiles below. The bounded short-window length \(L_2\) and threshold \(K=\lceil8(h+2)/\varepsilon\rceil\) give conditional failure
\[
 O\!\left(\delta^K+L\delta^{K+1}\right)=o(x^{-2}).
\]
Only \(L_2\) starts contain the anchor; each other start needs \(K+1\) retained sites. For a long window of size \(G'\), write \(k=N_1-1\). The source takes \(k\ge6eG'\delta\) and \(k\ge(h+3)\log_2x\). A union over \(k\)-subsets of the unconditioned sites gives a per-start bound \((1+o(1))(eG'A_*\delta/k)^k\), interpreted as zero if \(k>G'\). Since \(A_*/6<1/2\), the union over all starts still costs \(o(x^{-2})\). Summing the three failures against the anchor marginals and the original weights now yields \(\mathbb E B=o(M)\) whenever \(M>0\). Markov's inequality gives sets of probability tending to one on which \(B/M=o(1)\). Their cardinality term is \(O((\delta n)^{-1})\). In particular they pass the original private comparison \(S\le\mathbb E_\delta S\). This argument extends the comparison to profiles with zero-rate sites. The published independent-product bound assumes positive site rates in a prescribed interval and does not by itself establish this extension.

\subsection{What the gain pays on a missing edge}

Let \(b=\exp(x^\varepsilon)\), \(a=xb\), \(\lambda=b/H\), and \(b'=a/H=x\lambda\). With \(m\) lifts, put \(U=mNb\). On an active orbit \(v\) at a served high central prime and on one fixed lift \(j\), set
\[
 A=g_U-\mathrm{Pr},\qquad F=\max(g_0,0),\qquad R_{jvp}=(F-A)_+ .
\]
The credited local gain is \(\max(A,F)=A+R_{jvp}\). Introduce \(C_{jvp}\), the sum of all excluded positive demand columns together with the minimum columns at even sources. Splitting each odd-power column into its positive excess and its minimum gives \(g_U-g_0=C_{jvp}\), since the same floor loss \(dl\) occurs in both. Thus
\[
 R_{jvp}=\bigl(\mathrm{Pr}-C_{jvp}+(-g_0)_+\bigr)_+,
 \qquad 0\le(-g_0)_+\le dl .
\]
A larger \(C\) suppresses recovery. A missing edge matters because an immediate selected-to-selected successor is precisely a first visit at odd distance one. The definition of exclusion removes its positive excess from \(g_0\). An even source also contributes the minimum part to \(C\). When both pieces are present their sum is the full adjusted first supply.
\href{https://github.com/nousentllc/riemann-hypothesis/blob/40e56bab026552e089d7afcfeff08d2defb9c005/src/Formalism/SharpRestriction/GrowthBounds/IsResidueClassZero.lean#L228-L268}{Source definitions of exclusion, the two gains and the price}.
\href{https://github.com/nousentllc/riemann-hypothesis/blob/40e56bab026552e089d7afcfeff08d2defb9c005/src/Formalism/SharpRestriction/GrowthBounds/IsResidueClassZero.lean#L336-L365}{Local gain maximum and other ledger entries}.

Suppose neither oriented first edge is possible in the supported residues on \(v\). The first excluded positive excess vanishes: no retained source has a retained successor. At an even source the immediate target is thin and has cap \(O(\lambda/n+\eta+c_G)\), with tolerance \(\eta\) and shift \(c_G\) from the frame. All deeper odd powers together cost \(O(a/p^3)\) per retained source. For \(p\ge H/x\) this tail divided by \(\lambda\) is \(O(x^4/H^2)\). The scale window gives \((\eta+c_G)/\lambda=o(1)\). Consequently \(C_{jvp}=o(N_v\lambda)\) for retained count \(N_v\). If \(O_v\) denotes its odd count, the variable price is \(O_v\lambda+o(N_v\lambda)\). The fixed price and the floor loss also disappear after summing over active representatives at a given prime and lift. Indeed their length terms contain \(\sum_v L_v\le n\), and their retained terms contain \(\sum_v N_v\le N\); their sum is bounded by a constant times
\[
 \lambda+n(\eta+c_G+N_m)+nN_m^2\log_2(2a)+N\log_2^2(2a)=o(N\lambda),
\]
where \(N_m\) is polynomial in the scale. The first term comes from \(n(\lambda/n)\). Keeping it outside the per-retained-site estimate also treats very sparsely occupied orbits. This proves that the recovery of the blocked orbits equals their odd-count charge, up to an aggregate error \(o(N\lambda)\) per prime and lift.

\subsection{When the support permits an edge}

The converse needs a probabilistic step even when the residue computation is exact. If an orientation is allowed on a long orbit of length \(L\), periodicity supplies at least \(L/q\) disjoint potential pairs of that orientation. Indeed one admissible starting residue repeats at least once per \(q\) steps, with the same parity; in a homogeneous residue it suffices to take the appropriate half of the orbit. Success requires retaining \emph{both} endpoints. Under independent inclusions at rates \(k_t/(n/q)\), all indicators of these disjoint pairs are independent and each has probability at least \(c_1\delta^2\) for a fixed finite menu with positive minimum occupied rate. A multiplicative lower-tail bound, followed by the union over at most \(2n^2\) oriented step-orbit choices of length \(L\ge\sqrt{x}\), makes the probability of any such orientation having fewer than \(c_2 L\delta^2\) successful pairs at most
\[
 2n^2\exp(-c_3\sqrt{x}\,\delta^2).
\]
This estimate uses the independent law only for discovery of a likely event. Under that law each class count has integer mean \(k_t\) and mode \(k_t\), so imposing the exact target count costs a factor at most \(n+1\) per class. The fixed-count law consequently has failure probability tending to zero as well. Its edge event intersects the preceding score event. It is the \emph{same} retained mask that supplies all long-orbit edges and passes the original score.

An even-to-odd first hit pays at least \((a-2c_G)/H\ge b'/2\) eventually, since its excluded excess and minimum together supply the adjusted first column. When only the reverse orientation is available, the minimum need not be present. Its excluded excess is nevertheless at least \(b'/(2\sqrt{x})\) on the band
\[
 J_x=\left[\,H/x,\ H/(1+x^{-1/2})\,\right].
\]
In fact the supply is at least \(b'(1+x^{-1/2})(1-2c_G/a)\) there, and the cap at the retained target is at most \(b'(1+\beta_{\rm ht})+\lambda+\eta+c_G\), where the relative height tolerance \(\beta_{\rm ht}=O(x^{-2})\). This gives the claimed slack once \(x\) is large. The total contribution of the realized pairs, even at the smaller reverse-edge rate, is at least a constant times \((\sqrt{x}\delta^2)L\lambda\). The price is at most \((1+o(1))L\lambda\) and \(dl=o(L\lambda)\). This proves \(A>F\) on those served orbits on every lift there. At \(p\) arbitrarily close to \(H\), a reverse hit could supply little or no excess, which is precisely why the narrower upper endpoint matters.
\href{https://github.com/nousentllc/riemann-hypothesis/blob/40e56bab026552e089d7afcfeff08d2defb9c005/src/Formalism/SharpRestriction/GrowthBounds/LogarithmicSums/PartnerCountOfCentral.lean#L769-L828}{High-prime price and immediate columns}.
\href{https://github.com/nousentllc/riemann-hypothesis/blob/40e56bab026552e089d7afcfeff08d2defb9c005/src/Formalism/SizeWindowBrackets/WindowSettingSatisfiable.lean#L65-L117}{Underlying size windows}.

\subsection{Prime averaging, the unresolved witness, and the denominator}

Write \(\Delta=\mathrm{ASK}-\mathrm{Ga}\) for the signed global difference. The earlier two-sided transport compares it to the positive high-central charge \(\mathcal K_x\), preserving complementary orbits, tax and demand and adding back the recovery on the active interior only. With
\[
 I_x=[\,H/x,\ H/(1+2/x+\beta_{\rm ht})\,],\qquad
 \mathcal R_x=\frac{2}{U}\sum_{\substack{p\in I_x\\\text{served high central}}}
       \log p\sum_{\substack{v\text{ active}\\j\text{ a lift}}}R_{jvp},
\]
the comparison reads \(\Delta/U+\mathcal R_x=\mathcal K_x+o(1)\). The factors \(2\log p\) come from the rational-prime deficit ledger. The earlier eligibility test used \(\log(p)/p\). Interchanging them would change the inferential problem. In passing from the interior for the gap identity to \(J_x\) for the edge estimates one discards a strip of normalized log weight \(O((x^{-1/2}+H^{-1})\log H)=o(1)\), by counting all integers. The possibly much larger errors below \(\sqrt H\) in the signed gap are already controlled separately in that two-sided transport by \(O(\sqrt H\,b'\log(2a)/b)=o(1)\).
\href{https://github.com/nousentllc/riemann-hypothesis/blob/40e56bab026552e089d7afcfeff08d2defb9c005/src/Formalism/SharpRestriction/GrowthBounds/LogarithmicSums/NotSingularOfCentral.lean#L679-L728}{Summation over the local gain maximum}.
\href{https://github.com/nousentllc/riemann-hypothesis/blob/40e56bab026552e089d7afcfeff08d2defb9c005/src/Formalism/SizeWindowBrackets/ChainSlopeWindows/DiscriminantSettingSatisfied.lean#L162-L220}{Global residuals}.
\href{https://github.com/nousentllc/riemann-hypothesis/blob/3dedd2d81a54d046e3c12ef78df9d9fd0c2ad6dc/research/uriel/experiments/actual-gain-study/evidence/transport/TRANSPORT.md}{Prior written two-sided comparison}.

The prior pointwise prime-step comparison starts from the smoothed prime sums for the host's characters. Fourier inversion in the group with two quadratic signs and additive step gives the joint sign \((-1,-1)\) and each fixed step normalized log weight \(1/(4n)+O_D(x^{-D})\), for every fixed \(D\). Its coefficient follows because exactly one of the \(4n\) characters is principal. A three-point difference bound removes the smoothing, retaining the host-uniform error; terms at higher prime powers have normalized weight \(O(H^{-1/2}\log H)\). Since \(n\) grows only polynomially, \(D\) can absorb a sum over all steps, even when the directional test depends on the labels or selection. In a dyadic cyclic group, the number of steps of order below \(\sqrt{x}\) is \(O(\sqrt{x})\). Their normalized prime weight is thus \(o(1)\). The service and drop estimates are independently needed to remove omitted point-prime incidences of total log weight \(o(HN)\). Uniform prime weights on the full group alone say nothing about where the service omissions fall.
\href{https://github.com/nousentllc/riemann-hypothesis/blob/40e56bab026552e089d7afcfeff08d2defb9c005/src/Formalism/SharpRestriction/GlobalSchurBounds/Chosen.lean#L46-L51}{Host-wide analytic input}.
\href{https://github.com/nousentllc/riemann-hypothesis/blob/40e56bab026552e089d7afcfeff08d2defb9c005/src/Formalism/SharpRestriction/GlobalSchurBounds/Joints.lean#L855-L883}{Artin characters on units}.
\href{https://github.com/nousentllc/riemann-hypothesis/blob/ab18d856b6989f7bcb05e486ca720aa9016f3268/research/uriel/experiments/odd-position-study/evidence/FAMILY.md}{Prior directional and service comparison}.

At a quotient node \(v\), let \(w_v\) be the number of its leaves and \(b_v\) the sum of the rates in its odd child. Write \(d_v\) for the number of odd increments within that node's cyclic coordinate that have neither directed successor orientation. The node has step mass \(w_v/(2q)\) and its odd retained mass relative to \(N\) tends to \(b_v/q\). The blocked steps in it have mass \(d_v/q\). The subdivisions finer than the 64 classes are homogeneous under the sampling law. Conditional on a uniform point inside a given class, exactly half the sites are odd at every finer nontrivial level, for any labeling. Their total odd-exposure weight is \(1/(2q)-1/(2n)\). Sampling at fixed count gives variance \(O(N)\) for the sum of these bounded exposures; it therefore concentrates after division by \(N\) on the same high-probability intersection. There is no nontrivial blocked orientation inside a homogeneous positive class. The preceding local-charge calculation also subtracts already-paired immediate predecessors. Their average fraction is negligible, since \(\sum_s|T\cap(T+s)|=N^2\) for every finite set and \(N/n\to0\). Combining these facts and the prime factor \(2/4\) gives
\[
 \mathcal K_x=\frac{q+\sum_v w_v b_v}{4q^2}+o(1),
 \qquad
 \mathcal R_x=\frac{2\sum_v d_v b_v}{4q^2}+o(1).
\]
These formulae permit the labeling to vary with the host; the right sides can then vary too. In particular the recovery fraction is the quotient of the two displayed numerators up to \(o(1)\), since the tail supplies a uniform positive lower bound for the denominator. A certificate for all labelings becomes a lower bound on the \(\liminf\) of compatible selections. Without the simultaneous abundance proof, only the lower bound contributed by blocked directions would follow.

This deduction puts the experiment's boundary at a specific quantifier. For every sufficiently late suitable host, the positive-probability intersection supplies \emph{some} \(T\) admitting the estimates. The private averaging lemma can select any set below the uniform mean. The source returns neither the elements of that witness nor a minimizer of its score. Quotient computations determine which directed edges a profile forbids, rates for finite leaves, and a ratio; they cannot identify the actual \(TT\) of that existential call, an explicit first family point, or its measured prime ledger. The new kernel check in this study treats finite ordered-field accounting. The host-to-profile transfer above remains a written argument dependent on the cited arithmetic hypotheses.
